\chapter{The Three-Session Arc\\สามบทสนทนาเพื่อสร้างแรงส่ง}

ความสัมพันธ์ Mentoring ระยะยาวไม่อาจออกแบบทั้งหมดล่วงหน้า แต่สามครั้งแรกควรมีโครงสร้างพอที่จะสร้างความไว้วางใจ ระบุประเด็นสำคัญ และทำให้เกิดการลงมือทำ The Three-Session Arc ทำหน้าที่เป็นสะพานจากการรู้จักกันไปสู่ความสัมพันธ์ที่มีจุดหมายและจังหวะติดตามชัดเจน

\learningobjectives{
  \item ออกแบบสาม session แรกให้ต่อเนื่องกัน
  \item ใช้ success marker และงานระหว่าง session เพื่อสร้างแรงส่ง
  \item ปรับแผนตามความต้องการและความพร้อมของ Mentee
}

\section{ก่อนเริ่ม: สร้างข้อตกลง}
ก่อน Session 1 ควรแลกเปลี่ยน Profile หรือ Portfolio ฉบับปัจจุบัน พร้อมกำหนดวัตถุประสงค์เบื้องต้น ความถี่ ช่องทางติดต่อ การรักษาความลับ และขอบเขตบทบาท Mentor ควรแจ้งว่าคำถามบางข้ออาจท้าทาย และ Mentee มีสิทธิปฏิเสธหรือขอเปลี่ยนทิศทางได้

\section{Session 1: Listen First --- Build Trust}
\textbf{เป้าหมาย:} เข้าใจบุคคล บริบท ความหวัง และประเด็นจริงก่อนเสนอทางออก ให้ Mentee พูดเป็นสัดส่วนมากกว่า Mentor และจบด้วยประเด็นสำคัญ 1--2 เรื่อง ไม่ใช่รายการปัญหายาว

\begin{longtable}{P{.23\textwidth}P{.66\textwidth}}
\toprule
\textbf{องค์ประกอบ} & \textbf{แนวทาง}\\
\midrule
เวลา & 60 นาที\\
คำถามหลัก & \enquote{ส่วนใดของงานที่ท่านภูมิใจ และเพราะอะไร} \quad \enquote{อะไรยังไม่ชัดหรือทำให้ลังเล} \quad \enquote{ในหกเดือน ท่านอยากเห็นอะไรเปลี่ยน}\\
กิจกรรม & สำรวจ Tech Portfolio ทั้งเจ็ดแบบกว้าง และเลือก Gap ที่มีความหมาย\\
Success marker & Mentee กล้าพูดถึงประเด็นจริง และทั้งสองฝ่ายเข้าใจเป้าหมายร่วมกัน\\
ระหว่าง session & Mentee รวบรวม Evidence หรือเขียนคำถามที่ต้องการสำรวจต่อ\\
\bottomrule
\end{longtable}

\section{Session 2: Deep Dive --- Fill the Gap}
\textbf{เป้าหมาย:} เจาะ Gap สำคัญหนึ่งเรื่องด้วยหลักฐานและสร้างผลลัพธ์ที่จับต้องได้ เช่น Impact Statement, Evidence Map หรือ Process draft

\begin{longtable}{P{.23\textwidth}P{.66\textwidth}}
\toprule
\textbf{องค์ประกอบ} & \textbf{แนวทาง}\\
\midrule
เวลา & 75 นาที\\
คำถามหลัก & \enquote{หลังครั้งก่อน ท่านสังเกตอะไรเพิ่ม} \quad \enquote{เรารู้อะไรจริงและยังสมมติอะไร} \quad \enquote{หลักฐานใดจะเปลี่ยนการตัดสินใจ}\\
กิจกรรม & ใช้ GROW และทำงานบน artifact จริงร่วมกัน\\
Success marker & มี draft หรือการทดลองที่ Mentee เป็นเจ้าของ พร้อมเกณฑ์คุณภาพ\\
ระหว่าง session & Mentee ทดสอบ draft กับผู้เกี่ยวข้องหรือเก็บหลักฐานชิ้นถัดไป\\
\bottomrule
\end{longtable}

\section{Session 3: Design the Future --- Commit and Connect}
\textbf{เป้าหมาย:} สังเคราะห์บทเรียน กำหนด roadmap 3--6 เดือน และออกแบบรูปแบบความสัมพันธ์ต่อไป การจบสาม session ไม่จำเป็นต้องแก้ทุกเรื่อง แต่ควรทำให้ Mentee เห็นก้าวต่อไปและระบบสนับสนุน

\begin{longtable}{P{.23\textwidth}P{.66\textwidth}}
\toprule
\textbf{องค์ประกอบ} & \textbf{แนวทาง}\\
\midrule
เวลา & 60 นาที\\
คำถามหลัก & \enquote{อะไรเปลี่ยนไปจากครั้งแรก} \quad \enquote{milestone ใดสำคัญที่สุดใน 90 วัน} \quad \enquote{ใครต้องมีส่วนร่วม}\\
กิจกรรม & สร้าง roadmap ระบุ milestone, evidence, owner, risk และ next check-in\\
Success marker & Mentee สรุปทิศทางและก้าวแรกได้ด้วยตนเอง\\
หลัง session & ส่ง roadmap ฉบับตกลงและกำหนดรูปแบบติดตาม\\
\bottomrule
\end{longtable}

\section{บทสนทนาเมื่อแผนไม่เดิน}
\begin{examplebox}
\textbf{Mentor:} \enquote{เราเคยตกลงว่าจะทดสอบ draft ก่อนวันนี้ แต่ยังไม่เกิดขึ้น ผมไม่ต้องการรีบสรุปว่าเป็นเรื่องวินัย ท่านช่วยเล่าว่าเกิดอะไรขึ้นได้ไหม}\\
\textbf{Mentee:} \enquote{ฉันกังวลว่าข้อมูลยังไม่ดีพอจนไม่กล้าส่ง}\\
\textbf{Mentor:} \enquote{ดังนั้นคอขวดอาจไม่ใช่เวลา แต่เป็นเกณฑ์คำว่า ``ดีพอ'' เราจะออกแบบการทดสอบที่มีความเสี่ยงต่ำอย่างไร}\\
บทสนทนานี้รักษาความรับผิดชอบโดยไม่ใช้ความอับอายเป็นแรงผลัก
\end{examplebox}

\section{บันทึกที่เคารพความสัมพันธ์}
หลังแต่ละครั้ง บันทึกเฉพาะสิ่งจำเป็น: เป้าหมาย ประเด็นสำคัญ การตัดสินใจ action owner และวันติดตาม หลีกเลี่ยงการบันทึกการตีความส่วนบุคคลที่ไม่จำเป็น และตกลงกับ Mentee ว่าสิ่งใดแบ่งปันได้

\begin{mentornote}
ความต่อเนื่องเกิดจากการกลับมาถามถึงสิ่งที่ Mentee บอกว่าสำคัญ ไม่ใช่จากการกำหนดงานจำนวนมาก เลือก action ที่เล็กพอจะทำจริงและสำคัญพอจะสร้างการเรียนรู้
\end{mentornote}

\begin{keytakeaway}
สาม session แรกควรสร้างลำดับจากความไว้วางใจ สู่ความชัด และสู่การลงมือทำ โดยแต่ละครั้งมีผลลัพธ์ที่เชื่อมต่อกับครั้งถัดไป
\end{keytakeaway}

\begin{reflection}
ร่าง Three-Session Arc สำหรับ Mentee หนึ่งคน ระบุผลลัพธ์ของแต่ละครั้ง ความเชื่อมโยงระหว่างครั้ง และเงื่อนไขที่จะทำให้ต้องปรับแผน
\blankline\blankline
\end{reflection}

\chaptersummary{The Three-Session Arc ทำให้จุดเริ่มต้นของความสัมพันธ์มีทิศทางโดยไม่ทำให้แข็งตัว Session แรกฟังและสร้างความไว้วางใจ Session สองเจาะ Gap และสร้างงานจริง Session สามออกแบบอนาคตและระบบติดตาม}
